% Generated draft: requires mathematical and source review.
\documentclass[11pt]{article}
\usepackage[margin=25mm]{geometry}
\usepackage{fontspec}
\setmainfont{Latin Modern Roman}
\usepackage{amsmath,amssymb,mathtools,unicode-math}
\setmathfont{Latin Modern Math}
\usepackage{enumitem,needspace,xurl,hyperref}
\setlength{\parindent}{0pt}
\setlength{\parskip}{5pt}
\setlength{\emergencystretch}{3em}
\newcommand{\N}{\mathbb{N}}
\newcommand{\Z}{\mathbb{Z}}
\newcommand{\Q}{\mathbb{Q}}
\newcommand{\R}{\mathbb{R}}
\newcommand{\C}{\mathbb{C}}
\newcommand{\F}{\mathbb{F}}
\newcommand{\A}{\mathbb{A}}
\newcommand{\PP}{\mathbb{P}}
\newcommand{\E}{\mathbb{E}}
\newcommand{\eps}{\varepsilon}
\newcommand{\dd}{\,\mathrm{d}}
\begin{document}
\section*{UnsolvedMath ID: 30004812}
\textbf{Source problem code:} catalogue entry
\subsection{Definitions and mathematical statement}
The supplied record poses a research programme rather than a single fully specified conjecture. A knot in a three-manifold \(M\) is an embedded copy of \(S^1\). The source does not specify the category of embeddings, the equivalence relation on knots, whether knots or manifolds are oriented, or whether \(M\) is connected, compact, or allowed to have boundary. These choices must be fixed before a precise universal statement can be made.

For a specified class of three-manifolds and a specified knot equivalence relation \(\sim_M\), write
\[
\mathcal K(M)=\{K\subset M:K\text{ is an admissible embedded circle}\}/\!\sim_M.
\]
The source asks the following related questions for knots in general three-manifolds:
\begin{enumerate}
\item Can one construct Jones/HOMFLYPT-type polynomial invariants? A polynomial-valued invariant would be a map
\[
I_M:\mathcal K(M)\longrightarrow R_M,
\]
where \(R_M\) is a specified polynomial or Laurent polynomial ring. Neither this ring nor the properties that make an invariant Jones/HOMFLYPT-type are specified. In particular, merely exhibiting an arbitrary invariant would not resolve the intended question; its required normalization, skein relations, relationship with the classical polynomials, and dependence on auxiliary choices remain to be stated.
\item What projections of knots in \(M\), together with what diagrammatic information, have useful computational and invariant-theoretic properties? The projection target, admissible singularities, auxiliary structures, and meaning of usefulness are unspecified.
\item Is there a Reidemeister theorem for such projections? Once a diagram system and its allowed moves have been specified, the intended completeness criterion is that two diagrams represent equivalent knots if and only if they are related by the allowed diagram equivalences and moves. The source does not specify those moves or require a finite list, so no particular move theorem can be extracted from the record.
\end{enumerate}
These questions retain their unresolved qualifications; they do not assert the existence of a particular invariant or diagram calculus. The supplied dated curation labels the problem partially solved and reports established Jones-type generalizations through Kauffman-bracket and related skein modules, while reporting that a single canonical Jones/HOMFLYPT polynomial with the full classical role for arbitrary three-manifolds is not established. This is a supplied literature assessment, not an independently verified theorem or status determination.
\subsection{Short English statement}
The problem asks how to extend Jones/HOMFLYPT-type invariants and knot diagrams to general three-manifolds, and whether suitable diagram moves characterize knot equivalence. The manifold class, knot conventions, invariant requirements, projections, and moves remain unspecified. The supplied partially solved label is dated curation requiring independent review.
\subsection{Sources}
\begin{itemize}
\item Local UnsolvedMath registry, permanent ID \(30004812\), source problem code \texttt{catalogue entry}, title: Polynomial Knot Invariants and Projections in General Three-Manifolds. The supplied record labels the statement exact and reports byte-for-byte retention; this does not independently verify its attribution or mathematical precision. No individual proposer is supplied.
\item Supplied report citation: \emph{MATRIX-MFO Tandem Workshop: Invariants and Structures in Low-Dimensional Topology (2022)}, Oberwolfach Reports, pp.\ 2327--2364, DOI \texttt{10.4171/owr/2021/42}. \url{https://doi.org/10.4171/owr/2021/42}. The background describes Oberwolfach Reports as documenting workshops and their open-problem discussions. The registry separately gives proposed year 2021; both year indications are retained without reconciliation.
\item Supplied literature source labelled ``Detcherry'': \url{https://doi.org/10.1112/jlms.12410}. The record supplies no title, full author information, or specific supporting theorem.
\item Literature assessment checked at \texttt{2026-08-22}: supplied dated curation, not independently verified. Its listed sources are the OWR report and Detcherry links above. It assigns status partially solved, identifies Kauffman-bracket and related skein modules as established Jones-type generalizations, and states that a single canonical Jones/HOMFLYPT polynomial with the full classical role for arbitrary three-manifolds is not established. The background explicitly calls for independent verification of time-sensitive status and cited claims.
\end{itemize}
\end{document}
